%=========================================================
% C++ Chapter 7 Lab Handout (Verbatim Korean-safe)
% Topic: friend & operator overloading
% Compile with XeLaTeX
%=========================================================
\documentclass[11pt,a4paper]{article}

\usepackage[margin=1in]{geometry}
\usepackage{fontspec}
\usepackage{kotex}
\usepackage{hyperref}
\usepackage{enumitem}
\usepackage{fancyhdr}
\usepackage{titlesec}
\usepackage{xcolor}
\usepackage{tcolorbox}
\usepackage{fvextra}

% ---------- Fonts ----------
% If D2Coding isn't available, try: NanumGothicCoding, Noto Sans Mono CJK KR
\setmonofont{D2Coding}[Scale=MatchLowercase]

% ---------- Colors ----------
\definecolor{CodeBg}{RGB}{248,248,248}
\definecolor{CodeFrame}{RGB}{220,220,220}
\definecolor{Accent}{RGB}{40,80,160}

% ---------- Header ----------
\pagestyle{fancy}
\fancyhf{}
\lhead{\textbf{C++ 7장 실습}}
\rhead{\textbf{프렌드와 연산자 중복}}
\cfoot{\thepage}

\titleformat{\section}{\Large\bfseries\color{Accent}}{\thesection}{0.6em}{}
\titleformat{\subsection}{\large\bfseries}{\thesubsection}{0.6em}{}

% ---------- Boxes ----------
\newtcolorbox{GoalBox}{
  colback=CodeBg,
  colframe=CodeFrame,
  arc=3pt,
  boxrule=0.8pt,
  left=10pt,right=10pt,top=8pt,bottom=8pt
}
\newcommand{\filetag}[1]{\texttt{\color{Accent}#1}}

% ---------- Code block environment (Unicode-safe) ----------
\DefineVerbatimEnvironment{cppcode}{Verbatim}{
  fontsize=\small,
  frame=single,
  framerule=0.6pt,
  rulecolor=\color{CodeFrame},
  numbers=left,
  numbersep=8pt,
  baselinestretch=1,
  breaklines=true,
  breaksymbolleft=,
  breaksymbolright=,
  commandchars=\\\{\},
  xleftmargin=1.5em
}

\begin{document}

\begin{center}
  {\LARGE \textbf{C++ 7장 실습: 프렌드와 연산자 중복}}\\
  \vspace{6pt}
  {\large (friend function/class, operator overloading: member vs friend, binary/unary, prefix/postfix)}
\end{center}

\vspace{8pt}

\section*{학습 목표}
\begin{GoalBox}
\begin{enumerate}[leftmargin=18pt]
  \item 프렌드 함수의 개념을 이해하고, 프렌드 함수를 작성할 수 있다.
  \item 연산자 중복의 개념을 이해하고, 연산자를 중복할 수 있다.
  \item 연산자를 클래스 멤버 함수로 작성할 수 있다.
  \item 연산자를 프렌드 함수로 작성할 수 있다.
  \item 다양한 이항/단항 연산자를 중복할 수 있다.
  \item 단항 연산자에서 전위/후위를 구분하여 작성할 수 있다.
\end{enumerate}
\end{GoalBox}

\section*{실습 운영 규칙}
\begin{itemize}[leftmargin=18pt]
  \item 각 실습 코드는 \texttt{main()}을 포함하므로 독립 실행 가능합니다.
  \item 일부 코드는 ``왜 안 되는지'' 보여주기 위해 (에러 확인) 주석이 있습니다.
  \item 연산자 오버로딩은 \textbf{가독성 향상}이 목적입니다. 의미가 모호하면 오히려 해가 됩니다.
\end{itemize}

\tableofcontents
\newpage

%=========================================================
\section{프렌드 1: 외부 함수를 friend로 선언(예제 7-1)}
\begin{GoalBox}
\textbf{핵심}
\begin{itemize}[leftmargin=18pt]
  \item friend 함수는 클래스 멤버가 아닌 ``외부 함수''인데, private 멤버 접근 권한을 받습니다.
  \item 외부 함수가 클래스 타입을 인자로 받을 때, 클래스 정의 전에 타입을 알려주는 \textbf{전방 선언}(\texttt{class Rect;})이 필요할 수 있습니다.
\end{itemize}
\end{GoalBox}

\noindent\filetag{L1\_FriendGlobalEquals.cpp}
\begin{cppcode}
#include <iostream>
using namespace std;

class Rect;                       // (1) 전방 선언: Rect를 미리 "이름만" 알려줌
bool equals(Rect r, Rect s);      // (2) 외부 함수 원형

class Rect {
    int width, height;            // private(디폴트): 외부에서 접근 불가
public:
    Rect(int w, int h) : width(w), height(h) {}
    friend bool equals(Rect r, Rect s); // (3) equals를 friend로 선언 -> private 접근 허용
};

bool equals(Rect r, Rect s) {     // (4) 외부 함수 (멤버 함수 아님!)
    // friend라서 private인 width/height 접근 가능
    return (r.width == s.width && r.height == s.height);
}

int main() {
    Rect a(3,4), b(4,5);
    cout << (equals(a,b) ? "equal\n" : "not equal\n");
}
\end{cppcode}

%=========================================================
\section{프렌드 2: 다른 클래스의 특정 멤버 함수를 friend로(예제 7-2)}
\begin{GoalBox}
\textbf{핵심}
\begin{itemize}[leftmargin=18pt]
  \item ``RectManager::equals'' 같은 \textbf{특정 멤버 함수}만 friend로 초대할 수 있습니다.
  \item 문법: \texttt{friend bool RectManager::equals(Rect, Rect);}
\end{itemize}
\end{GoalBox}

\noindent\filetag{L2\_FriendOtherClassMember.cpp}
\begin{cppcode}
#include <iostream>
using namespace std;

class Rect; // 전방 선언

class RectManager {
public:
    bool equals(Rect r, Rect s);
};

class Rect {
    int width, height;
public:
    Rect(int w, int h) : width(w), height(h) {}
    friend bool RectManager::equals(Rect r, Rect s); // 특정 멤버만 friend
};

bool RectManager::equals(Rect r, Rect s) {
    return (r.width == s.width && r.height == s.height); // friend라서 접근 가능
}

int main() {
    Rect a(3,4), b(3,4);
    RectManager man;
    cout << (man.equals(a,b) ? "equal\n" : "not equal\n");
}
\end{cppcode}

%=========================================================
\section{프렌드 3: 다른 클래스 전체를 friend로(예제 7-3)}
\begin{GoalBox}
\textbf{핵심}
\begin{itemize}[leftmargin=18pt]
  \item \texttt{friend RectManager;} 형태로 \textbf{클래스 전체}를 friend로 줄 수 있습니다.
  \item 장점: RectManager의 여러 기능이 Rect의 private에 접근 가능
  \item 단점: 캡슐화가 약해지므로 남발 금지(최소 권한 원칙).
\end{itemize}
\end{GoalBox}

\noindent\filetag{L3\_FriendWholeClass.cpp}
\begin{cppcode}
#include <iostream>
using namespace std;

class Rect; // 전방 선언

class RectManager {
public:
    bool equals(Rect r, Rect s);
    void copy(Rect& dest, Rect& src);
};

class Rect {
    int width, height;
public:
    Rect(int w, int h) : width(w), height(h) {}
    friend class RectManager; // 클래스 전체 friend
};

bool RectManager::equals(Rect r, Rect s) {
    return (r.width == s.width && r.height == s.height);
}
void RectManager::copy(Rect& dest, Rect& src) {
    dest.width = src.width;
    dest.height = src.height;
}

int main() {
    Rect a(3,4), b(5,6);
    RectManager man;
    man.copy(b, a);
    cout << (man.equals(a,b) ? "equal\n" : "not equal\n");
}
\end{cppcode}

%=========================================================
\section{연산자 중복: 규칙 요약(꼭 알아야 하는 제한)}
\begin{GoalBox}
\textbf{중요 제한}
\begin{itemize}[leftmargin=18pt]
  \item C++에 원래 존재하는 연산자만 중복 가능(새 기호 생성 불가).
  \item 피연산자 개수 변경 불가(이항은 이항, 단항은 단항).
  \item 우선순위/결합법칙 변경 불가.
  \item 일부 연산자는 중복 불가(\texttt{.}, \texttt{::*}, \texttt{?:} 등).
\end{itemize}
\end{GoalBox}

%=========================================================
\section{연산자 함수 작성 1: 멤버 함수로 + 구현(예제 7-4)}
\begin{GoalBox}
\textbf{핵심}
\begin{itemize}[leftmargin=18pt]
  \item 멤버 함수로 이항 연산자 구현 시: \texttt{a + b} 는 \texttt{a.operator+(b)}로 변환됩니다.
  \item 오른쪽 피연산자만 매개변수로 받습니다(왼쪽은 \texttt{this}).
\end{itemize}
\end{GoalBox}

\noindent\filetag{L4\_PowerPlusMember.cpp}
\begin{cppcode}
#include <iostream>
using namespace std;

class Power {
    int kick;
    int punch;
public:
    Power(int k=0, int p=0) : kick(k), punch(p) {}
    void show() const { cout << "kick=" << kick << ",punch=" << punch << "\n"; }

    Power operator+(Power op2) const { // a + b  ==> a.operator+(b)
        Power tmp;
        tmp.kick = this->kick + op2.kick;
        tmp.punch = this->punch + op2.punch;
        return tmp;
    }
};

int main() {
    Power a(3,5), b(4,6), c;
    c = a + b;
    a.show(); b.show(); c.show();
}
\end{cppcode}

%=========================================================
\section{연산자 함수 작성 2: == 구현(예제 7-5)}
\begin{GoalBox}
\textbf{핵심}
\begin{itemize}[leftmargin=18pt]
  \item 비교 연산자(\texttt{==})는 보통 \textbf{const} 멤버 함수로 구현합니다.
  \item 성능: 인자를 값으로 받기보다 \texttt{const Power\&}가 일반적으로 더 좋습니다.
\end{itemize}
\end{GoalBox}

\noindent\filetag{L5\_PowerEqualMember.cpp}
\begin{cppcode}
#include <iostream>
using namespace std;

class Power {
    int kick, punch;
public:
    Power(int k=0, int p=0) : kick(k), punch(p) {}
    void show() const { cout << "kick=" << kick << ",punch=" << punch << "\n"; }

    bool operator==(const Power& op2) const { // a == b  ==> a.operator==(b)
        return (kick == op2.kick && punch == op2.punch);
    }
};

int main() {
    Power a(3,5), b(3,5);
    a.show(); b.show();
    cout << (a==b ? "same\n" : "diff\n");
}
\end{cppcode}

%=========================================================
\section{연산자 함수 작성 3: += 는 참조 리턴(예제 7-6)}
\begin{GoalBox}
\textbf{핵심}
\begin{itemize}[leftmargin=18pt]
  \item \texttt{a += b}는 \texttt{a} 자체를 바꾸므로 결과를 \textbf{\texttt{*this} 참조}로 리턴합니다.
  \item 그래서 \texttt{c = (a += b);} 같은 코드가 자연스럽게 동작합니다.
\end{itemize}
\end{GoalBox}

\noindent\filetag{L6\_PowerPlusEqual.cpp}
\begin{cppcode}
#include <iostream>
using namespace std;

class Power {
    int kick, punch;
public:
    Power(int k=0, int p=0) : kick(k), punch(p) {}
    void show() const { cout << "kick=" << kick << ",punch=" << punch << "\n"; }

    Power& operator+=(const Power& op2) { // a += b  ==> a.operator+=(b)
        kick += op2.kick;
        punch += op2.punch;
        return *this; // (핵심) 참조 리턴
    }
};

int main() {
    Power a(3,5), b(4,6), c;
    c = (a += b);
    a.show();
    c.show();
}
\end{cppcode}

%=========================================================
\section{실습: b = a + 2; (멤버 함수로 + 오버로딩)(예제 7-7)}
\begin{GoalBox}
\textbf{핵심}
\begin{itemize}[leftmargin=18pt]
  \item \texttt{a + 2}는 가능(왼쪽이 Power).
  \item 하지만 \texttt{2 + a}는 멤버 함수로는 불가능(왼쪽이 int라서 this가 될 수 없음) $\rightarrow$ 다음 섹션에서 friend로 해결.
\end{itemize}
\end{GoalBox}

\noindent\filetag{L7\_PowerPlusIntMember.cpp}
\begin{cppcode}
#include <iostream>
using namespace std;

class Power {
    int kick, punch;
public:
    Power(int k=0, int p=0) : kick(k), punch(p) {}
    void show() const { cout << "kick=" << kick << ",punch=" << punch << "\n"; }

    Power operator+(int op2) const { // a + 2  ==> a.operator+(2)
        return Power(kick + op2, punch + op2);
    }
};

int main() {
    Power a(3,5), b;
    b = a + 2;
    a.show();
    b.show();
}
\end{cppcode}

%=========================================================
\section{2 + a 를 가능하게: friend 외부 연산자 함수(예제 7-11)}
\begin{GoalBox}
\textbf{핵심}
\begin{itemize}[leftmargin=18pt]
  \item \texttt{2 + a}는 왼쪽이 \texttt{int}이므로 멤버로 구현 불가.
  \item 해결: \textbf{외부 함수} \texttt{operator+(int, Power)}를 만들고, Power에 friend로 선언해 private 접근 허용.
\end{itemize}
\end{GoalBox}

\noindent\filetag{L8\_PowerIntPlusFriend.cpp}
\begin{cppcode}
#include <iostream>
using namespace std;

class Power {
    int kick, punch;
public:
    Power(int k=0, int p=0) : kick(k), punch(p) {}
    void show() const { cout << "kick=" << kick << ",punch=" << punch << "\n"; }
    friend Power operator+(int op1, const Power& op2); // friend 선언
};

Power operator+(int op1, const Power& op2) { // 외부 함수
    return Power(op1 + op2.kick, op1 + op2.punch);
}

int main() {
    Power a(3,5), b;
    b = 2 + a;
    a.show();
    b.show();
}
\end{cppcode}

%=========================================================
\section{단항 연산자 1: 전위 ++ (멤버 함수)(예제 7-8)}
\begin{GoalBox}
\textbf{핵심}
\begin{itemize}[leftmargin=18pt]
  \item 전위: \texttt{++a} 는 \texttt{a.operator++()} 로 변환.
  \item 보통 \textbf{참조 리턴}(\texttt{Power\&})하여 체이닝을 가능하게 합니다.
\end{itemize}
\end{GoalBox}

\noindent\filetag{L9\_PrefixIncMember.cpp}
\begin{cppcode}
#include <iostream>
using namespace std;

class Power {
    int kick, punch;
public:
    Power(int k=0, int p=0) : kick(k), punch(p) {}
    void show() const { cout << "kick=" << kick << ",punch=" << punch << "\n"; }

    Power& operator++() { // ++a
        kick++; punch++;
        return *this;
    }
};

int main() {
    Power a(3,5), b;
    b = ++a;
    a.show();
    b.show();
}
\end{cppcode}

%=========================================================
\section{단항 연산자 2: ! (멤버 함수)(예제 7-9 실습)}
\begin{GoalBox}
\textbf{핵심}
\begin{itemize}[leftmargin=18pt]
  \item \texttt{!a}는 보통 bool을 리턴하도록 오버로딩합니다.
  \item 여기서는 kick,punch가 모두 0이면 true.
\end{itemize}
\end{GoalBox}

\noindent\filetag{L10\_NotOperatorMember.cpp}
\begin{cppcode}
#include <iostream>
using namespace std;

class Power {
    int kick, punch;
public:
    Power(int k=0, int p=0) : kick(k), punch(p) {}

    bool operator!() const { // !a
        return (kick == 0 && punch == 0);
    }
};

int main() {
    Power a(0,0), b(5,5);
    cout << (!a ? "a is zero\n" : "a not zero\n");
    cout << (!b ? "b is zero\n" : "b not zero\n");
}
\end{cppcode}

%=========================================================
\section{단항 연산자 3: 후위 ++ (멤버 함수)(예제 7-10)}
\begin{GoalBox}
\textbf{핵심}
\begin{itemize}[leftmargin=18pt]
  \item 후위: \texttt{a++} 는 \texttt{a.operator++(int)} 로 변환.
  \item 구분을 위해 \textbf{의미 없는 int 매개변수}를 둡니다.
  \item 후위는 ``증가 전 상태''를 값으로 리턴하는 것이 일반적입니다.
\end{itemize}
\end{GoalBox}

\noindent\filetag{L11\_PostfixIncMember.cpp}
\begin{cppcode}
#include <iostream>
using namespace std;

class Power {
    int kick, punch;
public:
    Power(int k=0, int p=0) : kick(k), punch(p) {}
    void show() const { cout << "kick=" << kick << ",punch=" << punch << "\n"; }

    Power operator++(int) { // a++
        Power tmp = *this;  // 증가 전 상태 저장
        kick++; punch++;
        return tmp;         // 증가 전 값 리턴
    }
};

int main() {
    Power a(3,5), b;
    b = a++;
    a.show(); // 증가 후
    b.show(); // 증가 전
}
\end{cppcode}

%=========================================================
\section{전위/후위를 friend로 구현(예제 7-13)}
\begin{GoalBox}
\textbf{핵심}
\begin{itemize}[leftmargin=18pt]
  \item friend로 구현 시, 왼쪽 피연산자도 매개변수로 받습니다.
  \item 전위: \texttt{operator++(Power\&)} $\rightarrow$ 참조 리턴
  \item 후위: \texttt{operator++(Power\&, int)} $\rightarrow$ 값 리턴(증가 전)
\end{itemize}
\end{GoalBox}

\noindent\filetag{L12\_IncFriend.cpp}
\begin{cppcode}
#include <iostream>
using namespace std;

class Power {
    int kick, punch;
public:
    Power(int k=0, int p=0) : kick(k), punch(p) {}
    void show() const { cout << "kick=" << kick << ",punch=" << punch << "\n"; }

    friend Power& operator++(Power& op);        // 전위
    friend Power operator++(Power& op, int);    // 후위
};

Power& operator++(Power& op) { // ++op
    op.kick++; op.punch++;
    return op;
}
Power operator++(Power& op, int) { // op++
    Power tmp = op;
    op.kick++; op.punch++;
    return tmp;
}

int main() {
    Power a(3,5), b;
    b = ++a; a.show(); b.show();
    b = a++; a.show(); b.show();
}
\end{cppcode}

%=========================================================
\section{참조를 리턴하는 << 연산자: 체이닝(예제 7-14)}
\begin{GoalBox}
\textbf{핵심}
\begin{itemize}[leftmargin=18pt]
  \item \texttt{a << 3 << 5 << 6;} 같은 체이닝을 만들려면 \textbf{참조 리턴}이 핵심.
  \item 여기서 \texttt{<<}는 출력 연산이 아니라, ``파워 증가'' 의미로 재정의한 예제입니다(가독성 주의!).
\end{itemize}
\end{GoalBox}

\noindent\filetag{L13\_ShiftLeftAddChain.cpp}
\begin{cppcode}
#include <iostream>
using namespace std;

class Power {
    int kick, punch;
public:
    Power(int k=0, int p=0) : kick(k), punch(p) {}
    void show() const { cout << "kick=" << kick << ",punch=" << punch << "\n"; }

    Power& operator<<(int n) { // a << n  (a를 바꾸고 a 참조 리턴)
        kick += n;
        punch += n;
        return *this;
    }
};

int main() {
    Power a(1,2);
    a << 3 << 5 << 6;
    a.show(); // kick=15, punch=16
}
\end{cppcode}

%=========================================================
\section*{마무리 체크리스트}
\begin{GoalBox}
\begin{itemize}[leftmargin=18pt]
  \item friend 선언 3종류(외부 함수 / 다른 클래스의 멤버 함수 / 다른 클래스 전체)를 설명할 수 있는가?
  \item 연산자 오버로딩 제한 4가지를 말할 수 있는가?
  \item 이항 연산자를 멤버/프렌드로 구현할 때 매개변수 구조가 어떻게 달라지는가?
  \item 전위++/후위++의 시그니처 차이를 설명할 수 있는가?
  \item += 와 \texttt{<<} 체이닝에서 왜 참조 리턴이 필요한가?
\end{itemize}
\end{GoalBox}

\end{document}
